\chapter{Préparation, Planification et Architecture Système}

\section{Introduction}
Le succès d'un projet d'intelligence artificielle pour la smart city repose autant sur la gestion, l'organisation et la fiabilité du code que sur la performance des modèles. Ce chapitre détaille la logique de structuration choisie, les arguments derrière les choix techniques, et comment l'ensemble du projet CIRCU·LI est planifié pour assurer robustesse, flexibilité et évolutivité. Il pose également le contexte global, la problématique, et la conduite du changement inhérente à l'adoption de ce système.

\section{Contexte et Motivation du Projet}
\subsection{Contexte du projet et ODD}
Notre projet CIRCU·LI s’inscrit dans le contexte des Objectifs de Développement Durable (ODD) et prend place au sein d’un environnement urbain et technologique en pleine mutation. Face aux défis croissants liés à la gestion de la mobilité et à l’optimisation des flux de circulation, notre motivation est d’explorer des approches innovantes pour améliorer la planification urbaine et la sécurité routière.

CIRCU·LI contribue directement aux Objectifs de Développement Durable (ODD) suivants :
\begin{itemize}
    \item \textbf{ODD 9 (Industrie, Innovation et Infrastructure)} : CIRCU·LI introduit des solutions innovantes de gestion des flux urbains via la collecte en temps réel et l’analyse intelligente.
    \item \textbf{ODD 11 (Villes et Communautés durables)} : En réduisant la congestion, le projet contribue à rendre les villes plus sûres et plus accessibles.
    \item \textbf{ODD 13 (Lutte contre le Changement Climatique)} : L'optimisation des flux diminue les émissions de gaz à effet de serre.
    \item \textbf{ODD 8 (Travail décent et croissance économique)} : La réduction de la pollution liée au trafic participe au bien-être des citoyens (environnement sain, réduction du stress).
    \item \textbf{ODD 17 (Partenariats)} : Le projet met en œuvre des collaborations entre municipalités, opérateurs et acteurs technologiques.
\end{itemize}

\subsection{Problématique}
La gestion de la mobilité urbaine constitue aujourd’hui l’un des défis majeurs des collectivités locales. Avec l’augmentation continue du nombre de véhicules et l’urbanisation croissante, les villes font face à une congestion permanente. Malgré la présence de caméras, ces données restent sous-exploitées. Cette absence d’un système centralisé entraîne une réaction tardive aux problèmes de circulation et l'insatisfaction croissante des citoyens. CIRCU·LI naît de la nécessité de transformer ces données brutes en informations exploitables pour une gouvernance proactive.

\subsection{Motivation}
La croissance rapide du trafic exige des solutions plus intelligentes. Notre motivation provient de la volonté d'exploiter pleinement les données de mobilité pour rendre la ville fluide. CIRCU·LI est motivé par la vision de faire évoluer la gestion de la mobilité d’un modèle réactif vers un modèle prédictif et proactif.

\section{Conduite du changement et Stratégie (BSC)}
\subsection{Conduite du changement}
L’introduction d’une plateforme intelligente nécessite une conduite du changement structurée.
\begin{enumerate}
    \item \textbf{Identification des acteurs} : Municipalités, opérateurs de mobilité, analystes et citoyens.
    \item \textbf{Sensibilisation} : Présentation des objectifs et bénéfices.
    \item \textbf{Formation} : Initiation des décideurs aux indicateurs BI.
    \item \textbf{Mise en œuvre progressive} : Phase pilote sur un périmètre limité.
    \item \textbf{Évaluation} : Suivi de l'adoption et de l'impact sur la réduction de la congestion.
\end{enumerate}

\subsection{Alignement stratégique (Balanced Scorecard - BSC)}
L’approche Balanced Scorecard permet de traduire la vision stratégique en objectifs mesurables :
\begin{itemize}
    \item \textbf{Financier} : Optimiser les coûts d'exploitation IT (cible : réduction de 20\% en 2 ans).
    \item \textbf{Clients} : Fournir des tableaux de bord clairs (cible : satisfaction > 85\%).
    \item \textbf{Processus internes} : Automatiser l'analyse vidéo (cible : Uptime $\geq$ 99\%).
    \item \textbf{Apprentissage} : Développer une culture d'innovation via la formation IA/BI.
\end{itemize}

\section{Étude comparative et technologique}
\subsection{Benchmarking}
L’objectif de ce benchmarking est d’analyser des solutions existantes pour positionner CIRCU·LI. Par exemple, des solutions comme Citilog (France) ou Miovision (Canada) sont très coûteuses et peu adaptables au contexte local. CIRCU·LI se démarque par une utilisation des caméras existantes, une forte intégration BI et un coût maîtrisé.

\subsection{Technologie et outils}
La conception repose sur une sélection rigoureuse : Python, OpenCV et \textbf{YOLOv8} pour l'IA ; FastApi pour les Web Services ; PostgreSQL pour le stockage ; \textbf{Power BI / Grafana} pour la visualisation ; et Docker pour la conteneurisation. Le choix d'un serveur local Ubuntu garantit un traitement vidéo à faible latence (Edge Computing).

\section{Gestion de Projet Agile : Hybridation Scrum et XP}
Pour répondre à l'incertitude inhérente aux projets d'intelligence artificielle, nous avons opté pour une approche hybride combinant Scrum et Extreme Programming (XP). Le développement s'est réparti en \textbf{4 sprints} successifs.

\subsection{Planification macroscopique}
\begin{table}[H]
    \centering
    \begin{tabular}{|c|p{8cm}|l|}
        \hline
        \textbf{Sprint} & \textbf{Objectif} & \textbf{État} \\
        \hline
        1 & Benchmarking des algorithmes (Étude, évaluation et choix de YOLOv8) & Terminé \\
        2 & Socle technique et sécurité (Déploiement Ubuntu, Authentification, Rôles) & Terminé \\
        3 & Intelligence Artificielle et BI (Dataset, Entraînement, Dashboards) & Terminé \\
        4 & Intégration et Recommandations (API IA, Moteur de recommandations) & En cours \\
        \hline
    \end{tabular}
    \caption{Découpage en 4 Sprints du projet Circuli}
\end{table}

\subsection{Backlog du produit (Détail)}
\begin{table}[H]
    \centering
    \resizebox{\textwidth}{!}{
    \begin{tabular}{lllp{6cm}lc}
        \toprule
        \textbf{ID} & \textbf{EN TANT QUE…} & \textbf{JE SOUHAITE…} & \textbf{DANS LE BUT DE…} & \textbf{PRIORITÉ} & \textbf{Points} \\
        \midrule
        B01 & Chercheur & évaluer les algorithmes existants & choisir le modèle le plus adapté & Élevée & 5 \\
        S01 & Système & authentifier les utilisateurs & sécuriser l'accès à la plateforme & Élevée & 5 \\
        S02 & Administrateur & gérer les comptes et rôles & contrôler les accès & Élevé & 5 \\
        P01 & Analyste & collecter et structurer les données & alimenter Power BI / Grafana & Élevé & 3 \\
        P02 & Analyste & créer des tableaux de bord & visualiser les indicateurs clés & Moyenne & 5 \\
        IA01 & Système & détecter auto. les types de véhicules & reconnaître les véhicules & Élevé & 8 \\
        IA02 & Data Scientist & entraîner et valider le modèle & garantir la précision & Élevé & 8 \\
        IA03 & Développeur & intégrer le module IA & utiliser la reconnaissance & Moyenne & 5 \\
        R01 & Resp. mobilité & recevoir des recommandations & optimiser les déplacements & Élevé & 5 \\
        \bottomrule
    \end{tabular}
    }
    \caption{Product Backlog détaillé}
\end{table}

\subsection{Planification détaillée des Sprints}

\subsubsection*{Sprint 1 : Benchmarking des algorithmes}
\begin{table}[H]
    \centering
    \resizebox{\textwidth}{!}{
    \begin{tabular}{llc}
        \toprule
        \textbf{USER STORY (ID)} & \textbf{TÂCHES} & \textbf{DURÉE (JOURS)} \\
        \midrule
        \textbf{US-B01 – Benchmarking} & Recherche bibliographique (État de l'art) & 3 \\
        & Comparaison des modèles (YOLO, SSD, Faster R-CNN) & 4 \\
        & Évaluation des contraintes temps réel (FPS vs Précision) & 2 \\
        & Sélection et justification du modèle YOLOv8 & 1 \\
        \bottomrule
    \end{tabular}
    }
    \caption{Détail des tâches du Sprint 1}
\end{table}

\subsubsection*{Sprint 2 : Socle technique et sécurité (Ubuntu)}
\begin{table}[H]
    \centering
    \resizebox{\textwidth}{!}{
    \begin{tabular}{llc}
        \toprule
        \textbf{USER STORY (ID)} & \textbf{TÂCHES} & \textbf{DURÉE (JOURS)} \\
        \midrule
        \textbf{US-S01 – Authentification} & Configuration serveur Ubuntu et base de données & 2 \\
        & Développement login \& logout & 2 \\
        & Gestion sessions et tokens & 1 \\
        \midrule
        \textbf{US-S02 – Comptes \& rôles} & Définition rôles et permissions & 1 \\
        & Développement CRUD utilisateurs \& Interface admin & 3 \\
        \midrule
        \textbf{US-P01 – Collecte données} & Identification sources de données vidéo & 1 \\
        \bottomrule
    \end{tabular}
    }
    \caption{Détail des tâches du Sprint 2}
\end{table}

\subsubsection*{Sprint 3 : Intelligence Artificielle et BI}
\begin{table}[H]
    \centering
    \resizebox{\textwidth}{!}{
    \begin{tabular}{llc}
        \toprule
        \textbf{USER STORY (ID)} & \textbf{TÂCHES} & \textbf{DURÉE (JOURS)} \\
        \midrule
        \textbf{US-P02 – Dashboards BI} & Connexion base de données \& Création dashboards Grafana & 2 \\
        \midrule
        \textbf{US-IA01 – Détection} & Constitution dataset \& Annotation & 3 \\
        & Entraînement initial sur serveur Ubuntu & 3 \\
        \midrule
        \textbf{US-IA02 – Validation} & Optimisation hyperparamètres \& Évaluation & 4 \\
        \bottomrule
    \end{tabular}
    }
    \caption{Détail des tâches du Sprint 3}
\end{table}

\subsubsection*{Sprint 4 : Intégration et Recommandations}
\begin{table}[H]
    \centering
    \resizebox{\textwidth}{!}{
    \begin{tabular}{llc}
        \toprule
        \textbf{USER STORY (ID)} & \textbf{TÂCHES} & \textbf{DURÉE (JOURS)} \\
        \midrule
        \textbf{US-IA03 – Intégration IA} & Développement API IA (FastAPI) & 3 \\
        & Tests d'intégration & 1 \\
        \midrule
        \textbf{US-R01 – Recommandations} & Analyse besoins \& Moteur de recommandation & 3 \\
        & Visualisation \& Tests utilisateurs & 2 \\
        \midrule
        \textbf{US-AM – Amélioration} & Corrections, Optimisation Docker \& Documentation & 4 \\
        \bottomrule
    \end{tabular}
    }
    \caption{Détail des tâches du Sprint 4}
\end{table}

\section{Diagrammes UML et modélisation}
Pour formaliser les interactions entre les acteurs et la plateforme CIRCU·LI, nous avons modélisé les processus métiers via le langage UML.
\subsection{Diagramme de cas d'utilisation}
\begin{figure}[H]
  \centering
  \includegraphics[width=0.75\textwidth]{media/usecase_umldiag.png}
  \caption{Diagramme de cas d'utilisation global de la plateforme Circuli}
\end{figure}

\subsection{Diagramme de classes}
\begin{figure}[H]
  \centering
  \includegraphics[width=0.75\textwidth]{media/class_diag.png}
  \caption{Diagramme de classes UML de l'architecture logicielle}
\end{figure}

\section{Conclusion}
La combinaison d'une étude de contexte rigoureuse (ODD, conduite du changement) avec une architecture agile documentée assure à CIRCU·LI sa viabilité. Du choix des technologies (YOLO, BI) jusqu'à l'hybridation Scrum/XP répartie sur 4 Sprints, tout a été orchestré pour garantir une livraison robuste. Les diagrammes UML consolident l'approche qui sera implémentée dans le chapitre suivant.
